\subsection{A Last-Bar Same-Day Option Reading}
\label{sec:close_option_methods}

The one-bar forecast can be read against a listed price on the last
thirty minutes of an expiration day. What follows fixes the instrument,
the two measurements, the position rules, and the
control; constants were set before the portfolio was scored.

\paragraph{Instrument.} S\&P~500 weekly (SPXW) options with zero days
to expiry, cash-settled on the day's official close. The panel is the
expiration-day chain at 15:30~ET (bid, ask, mid, underlying, quoted
implied volatility), 2020 through 2025. At 15:30 the trade is the
nearest out-of-the-money pair: the smallest listed strike
$K_c \ge S$ with a live mid and the largest listed strike
$K_p \le S$ with a live mid, where $S$ is the 15:30 underlying print.
When $S$ sits on a listed strike the two coincide (a straddle);
otherwise the package is a one-strike-wide strangle. Entry
$P_{15:30}$ is the sum of the two mids. The position is held to the
close and settled at intrinsic against the official S\&P~500 close
$S_{\mathrm{cl}}$,
\begin{equation}
    \Pi_{\mathrm{long}}
    \;=\;
    \bigl(S_{\mathrm{cl}}-K_c\bigr)^{+}
    + \bigl(K_p-S_{\mathrm{cl}}\bigr)^{+}
    \;-\; P_{15:30},
    \label{eq:close_opt_settle}
\end{equation}
with no delta hedge and no share unwinding. The long return is
$R = \Pi_{\mathrm{long}}/P_{15:30}$; a short earns $-R$. By
construction $R \ge -1$, with equality on the days the close settles
between $K_p$ and $K_c$, either endpoint included, and both legs expire
worthless. The
16:00 chain underlying is retained only as a check on the settlement
print. Early-close sessions are dropped. The session ends at 13:00 on those
days, so the 15:30 chain row is stamped after the close and carries stale
quotes: the settlement is already known there and the short is a certain
$+1$. The test is a non-positive time to expiry on the 15:30 row --- the
vendor's \texttt{early\_close} column is true on every row of this file and
is not used. Twelve such sessions appear in the chain.
Six of them fall inside the forecast window, and the
five of those that carry a $16{:}00$ forecast row are the days
the filter takes out of the scored frame; the remaining one never had a
forecast row to lose.

\paragraph{Forecast remaining variance.} Eight per-bar forecast series
enter the same rules on the same contracts. Seven are estimators the body
of the paper defines: the baseline (HAR + calendar OLS) forecast
(Section~\ref{sec:dp_benchmark}), the block-diagonal ridge
(HAR versus exogenous block, separate penalties;
Section~\ref{sec:block_tikhonov}), LightGBM and XGBoost on the wide
all-features design (Section~\ref{sec:tree_edge}, frozen expert menu),
a causally tuned lasso on that same design under the selection
protocol of the paper's tuned ridge (Section~\ref{sec:dense_weak}), a
lasso at a fixed $\alpha = 10^{-4}$, and a causally tuned elastic net
--- the interpolating $\ell_1/\ell_2$ family --- on the same wide
all-features design (battery grid \texttt{logspace(-6,-2,5)} at
\texttt{l1\_ratio=0.5}, arm \texttt{b3\_enet\_tuned}). The eighth is not
a competing model: it is the block-diagonal ridge again, refit on a design
that carries no FOMC calendar columns, and it is carried as a diagnostic
row --- ``block-diagonal ridge, without the FOMC columns'' --- because the
gap between it and the ridge of record is what those calendar channels are
worth on this trade, and nothing else in the table isolates them. All eight
are stored on the estimation scale $y=\sqrt{RV/B}$. The 15:30 row is mapped
to raw thirty-minute variance by a causal second-order back-transform
\emph{adapted from} the evaluation protocol of
Section~\ref{sec:smearing_contract}: a Mincer--Zarnowitz map
$m=a+b\,\widehat{y}$ and a residual second moment $\hat\sigma^2$, both
fit on the previous 250 trading sessions only, then
$\widehat{RV}=(m^2+\hat\sigma^2)B$, an estimate of $E[RV]$ for the
15:30--16:00 bar. Adapted, not inherited, and the three differences are
stated rather than buried. The protocol's map is an ordinary least-squares
fit over the preceding \emph{evaluation window}, and its second moment
opens at that window's mean squared residual and then decays
exponentially within the window at $\lambda = 0.6$
(Equations~\eqref{eq:mean_calib}--\eqref{eq:var_ewma}). The map used here
is a \emph{weighted} fit, over a flat lookback of $250$ trading
\emph{sessions}, and its second moment is that same window's residual
variance held fixed across it. The differences are what a daily exercise
can refit --- this section scores 866 daily observations, not a quarter
million bars --- but they mean no number in this section is produced by the
map that produced the QLIKE tables, and none is carried between the two.
Forecast files end on 30~April~2024. The chain itself runs to
31~December~2025, so the 413 expiration days from 1~May~2024
onward --- the most recent twenty months of contiguous tape --- carry no
forecast row and are not scored anywhere below.

Two conventions matter for honesty. \emph{Alignment}: forecast rows
are labelled by the end of the bar they predict, so the row read at
15:30 is the one stamped 16:00 --- issued at 15:30 for the thirty
minutes that remain --- and not the stamp-15:30 row, which was issued
at 15:00 for a bar already finished; every input to the signal below
is therefore known at 15:30. \emph{Weighting}: the map and the
residual moment solve weighted least squares with weights
$1/\max(\widehat{y}_u,\,q_{10})^{2}$, where $q_{10}$ is the tenth
percentile of $\widehat{y}$ within the same trailing window, and the
fit window is restricted to regular-session bars. \emph{Regular session}
is an explicit mask, not an assumption, because the panel underneath is a
futures panel (Section~\ref{sec:dp_data}): a bar enters the fit only if
its own stamp lies in $10{:}30$--$16{:}00$, if its date carries a
$16{:}00$ stamp at all --- which drops cash-market holidays, whose futures
session ends at $13{:}00$ or $11{:}30$ --- and if its date is not one of
the NYSE $13{:}00$ early closes of 2001--2025, which drops the post-close
bars those days print through $16{:}00$. The mask is not cosmetic. Of the
$69{,}298$ bars the stamp window alone would admit it removes $1{,}083$,
which is $1.56\%$ of them: $510$ rows on $48$ of the $55$ early-close dates
in the 2001--2025 calendar, and $573$ rows on a further $125$ dates that
carry no $16{:}00$ stamp at all --- exchange holidays, whose rows here are
futures-only bars, and a few dates whose tape stops early. Those bars carry
far less variance than a regular one, so in a weighted fit they enter at a
low $y$ and drag the calibrated forecast down on every scored row.

One related contamination is disclosed rather than repaired. The per-slot
factor $B$ that scales the target is a trailing twenty-session mean of the
same slot, and it is built upstream of this section, before any of the
masking above: on an early close it absorbs that day's post-close bars, so
for the following twenty sessions the $16{:}00$ slot's $B$ sits a few
percent low. Ninety-four of the 866 scored days fall inside such a window.
The effect on the trade is second order --- $\widehat{RV}=(m^2+\hat\sigma^2)B$
moves with $B$, but so does the $y$ the map is calibrated against, and the
two largely cancel --- and a separate audit bounded what survives the
cancellation at under a tenth of a Sharpe ratio and a handful of flipped
positions. Repairing it at the source means regenerating the forecast
tables, and it is part of the pending campaign.

The weighting
addresses a level effect --- an unweighted fit is dominated by the
largest bars --- and was adopted on a forecast-accuracy criterion, not on
any portfolio outcome. A median-targeting map (the same line fitted by weighted least
absolute deviations, $\widehat{RV}=m^2 B$) was tested and not adopted:
it is right about the sign more often but earns less, because the
package's payoff is carried by the right tail of realized variance
that a median discards, and it cuts the $\mathrm{sign}(s)$ portfolio's
Sharpe ratio for all eight columns (block-diagonal ridge 1.34 to
0.49), though no one of those differences is resolved at 95\% by a
circular-block bootstrap.

\paragraph{Which panel produced which row.} The eight columns were not
all produced on one design, and the tables say which each came from. The
\emph{FOMC panel} --- the $1{,}264$-column design that carries the FOMC
calendar block --- is the panel of record here: the block-diagonal ridge
and the fixed lasso are fitted on it. The causally tuned lasso, the
elastic net, LightGBM and XGBoost are still on the $1{,}144$-column
\emph{earlier} panel, which has no such block; a run of those four on the
panel of record is pending, and the tables mark them so. The baseline is
invariant to the swap --- it reads no exogenous column, and re-derived on
both panels it agrees to $2\times 10^{-11}$ --- so its row is unaffected.
For an arm that does read exogenous columns the swap is not cosmetic: on
the ridge it moves the $16{:}00$ forecast by a median of $1.1\%$, and by
more than $1\%$ on $53\%$ of trading days. That is why the diagnostic row
exists, and why a ranking that mixes rows from the two panels would not be
a like-for-like ranking. Each table's file, panel, chunk tree, code
fingerprint and refit contract is recorded in
\texttt{results/spxw\_pnl/MANIFEST.md}.

\paragraph{Refit contract, hyperparameters, and what is not recoverable.}
Every linear arm is refit at each bar on a rolling window of $24{,}000$
bars (about $480$ trading sessions); the tuned lasso and elastic net
re-select their penalty from the fixed grid every $250$ solves, on the
embargoed forward tail of the current window. The two tree columns are
reported as shipped and are less clean than that. Their settings are
indices $00$ and $16$ into a frozen twenty-arm menu whose file is absent
both from the working tree and from the repository's history, so the
hyperparameters behind those two columns cannot be recovered; the two arms
were hand-picked from that bank and the basis of the pick is not recorded.
Their refit cadence is mixed as well: one cluster refit at every bar and
another every tenth bar, both contributed chunks to the same harvest, and
neither the chunk metadata nor the shipped table records which cadence
produced which stretch. A cadence coarser than one bar is stale but still
causal --- the fit still sees only strictly prior data, and an instrumented
replay found no look-ahead at either cadence --- so this is a provenance
defect and not a leak. One caveat attaches to the ridge of record too: its
table came from a cluster run identified by code fingerprint rather than by
a commit. The same estimator specification reproduces the corresponding
earlier-panel chunks to $2\times 10^{-13}$, so the estimator is confirmed;
what cannot be re-derived on a workstation is the panel build under that
exact code. A campaign that regenerates all eight columns under one
committed revision --- tree menu re-frozen, and cadence, thread count,
library versions and code hash written into every chunk --- is the stated
remedy, and it has not been run.

Two facts about the FOMC block itself belong with that. The release feed
behind it ends on 1~November~2023, so 122 of the 866 scored
expiration days --- $14.1\%$ --- fall after the feed is dead, and the
meetings of
13~December~2023, 31~January~2024 and 20~March~2024 carry no FOMC signal at
all. And the release-\emph{surprise} channel is a dead constant over the
whole panel: it is logged and then winsorized on a rolling $240$-bar
window, which clips a single-bar event to the running bound every time. So
``FOMC in the design'' means, operationally, the meeting-day indicator and
the bars-until and bars-since counters, and it means them only through
1~November~2023. The diagnostic row above measures the calendar channels as
they actually are in the shipped panel, not as they would be with a feed
that reaches the end of the sample.

\paragraph{What the 16:00 target is.} The forecast series predict $RV$ on
the paper's own panel, which is built from a near-24-hour futures tape
(Section~\ref{sec:dp_data}). The row read at 15:30 is therefore the
realized variance of the futures $15{:}30$--$16{:}00$ bar; it is not a
measurement of the cash market's settlement procedure, and it is nonzero on
days the cash market is shut. The option, by contrast, settles at the
official S\&P~500 close. On a regular session the two describe the same
half hour of the same index and the comparison is the intended one; the fit
mask above exists precisely to keep the sessions on which they are not
aligned out of both the calibration and the trade.

\paragraph{Quoted implied volatility.} Implied volatility is the
chain's quoted value on each of the two legs at 15:30; the package
quote is the mean of those two. The vendor quote is an \emph{hourly}
volatility while the window that remains is thirty minutes, so the
quote is rescaled before it is compared with anything,
\begin{equation}
    \mathrm{IV}_{30,t}
    \;=\;
    \mathrm{IV}^{\mathrm{hourly}}_t \big/ \sqrt{2},
    \label{eq:close_opt_ivscale}
\end{equation}
and $\mathrm{IV}_{30,t}^{2}$ is the implied variance of the remaining
bar --- the same window, and the same units, as $\widehat{RV}_t$. The
vendor field is not always a quote: it is the output of a bisection on
$[0.0005,\,0.025]$ that returns a bracket node when it fails to converge, so
the two bounds and the first interior nodes of that bisection are censored
values rather than prices, and they are masked here. On the five scored
days where a nearest leg carries one, the implied volatility is recovered
instead by inverting the Black--76 price of the package --- call plus put at
zero rate on the 15:30 index level, the vendor's own convention, over the
remaining half hour --- for the total volatility that reproduces the 15:30
midpoint. On uncensored legs that same inversion
reproduces the quoted midpoint to quote precision (median model-to-quote
price ratio 1.0000, 5th to 95th percentile 0.9969 to
1.0062), so ``quoted'' and ``inverted'' implied volatility coincide
there. Dropping the censored days outright instead of recovering them moves
no $\mathrm{sign}(s)$ Sharpe ratio in
Section~\ref{sec:close_option_results} by more than 0.06.

\paragraph{Signal.} The two measurements are compared in variance
space,
\begin{equation}
    s_t \;=\; \widehat{RV}_t - \mathrm{IV}_{30,t}^{2},
    \qquad
    \mathrm{VRP}_t \;=\; \mathrm{IV}_{30,t}^{2} - \widehat{RV}_t
    \;=\; -\,s_t,
    \label{eq:close_opt_signal}
\end{equation}
so $s_t>0$ says the forecast is above the price of the remaining
variance and $s_t<0$ says the premium is positive. The daily portfolio
return is $R'_t = q_t R_t$ for a position $q_t$ fixed at 15:30 from
information dated $t$ or earlier.

\paragraph{Position rules.} Two rules are scored on the same $R_t$ and
the same days; only $q_t$ changes.
\begin{itemize}
\itemsep2pt
\item \emph{Always short}: $q_t = -1$ every day. The rule reads no
      forecast --- it is the variance-risk-premium portfolio with nothing
      in it --- and it is the control for the other. It is called the
      always-short control throughout, and it is one row of every table
      because it is the same portfolio whichever forecast sits beside it.
\item \emph{$\mathrm{sign}(s)$}: $q_t = +1$ when $s_t > 0$ and
      $q_t = -1$ otherwise --- long the package when the forecast exceeds
      implied variance, short when it does not, with the tie at $s_t = 0$
      resolved short so that the position is never flat. Direction only, at
      unit size. This is the headline rule.
\end{itemize}
% VOL-TARGET PARKED 2026-09-02: overlay held out of the deck and the paper by author order.
\iffalse
\paragraph{Sizing overlay.} A portfolio that holds $|q_t| = 1$ every day
inherits the market's volatility cycle. One overlay is scored on top
of each rule, and it reads no forecast: let $\hat\sigma_t$ be the
standard deviation of the rule's own daily return over the trailing
63 sessions, lagged one day. The position is scaled by
\begin{equation}
    \ell_t \;=\;
    \min\!\left(
        \frac{\operatorname{med}_{u \le t-1} \hat\sigma_u}{\hat\sigma_t},
        \, 3
    \right),
    \label{eq:close_opt_voltarget}
\end{equation}
the expanding median of $\hat\sigma$ over its own past divided by its
current value. Average leverage is near one by construction and there
is no free target parameter; the cap at three never binds in sample.
The first sessions before an estimate exists stay flat, and scaled and
unscaled portfolios are compared on the same remaining days. Because the
overlay claims no forecast information, it is judged on risk
stability --- the wander of the portfolio's rolling volatility, its
drawdown, its per-year volatility --- and not on the Sharpe ratio,
which a pure rescaling should leave roughly unchanged.
\fi

% IRON-FLIES PARKED 2026-09-03: cost at every width, fraction of wealth unchanged; deck section parked, lab remains in the experimental notebook.
\iffalse
\paragraph{Defined-risk variant.} A retail account trading cash-settled
SPX or XSP under defined-risk margin cannot hold the short package
naked: every short leg must be covered. The variant scored here is that
constraint made explicit. On selling days the rule buys a wing on each
side of the body --- the nearest strike with a live 15:30 mid at least
$w$ points further out of the money, $w \in \{20, 25, 30, 50\}$ --- so
the position is an iron butterfly around the same package as everywhere
else in this section (an iron condor one strike wide when the body is a
strangle). The net credit is $C$, body premium minus wing premium, and
days with $C \le 0$ are dropped with a printed count. The worst
settlement loses the larger \emph{actual} wing gap minus the credit ---
the nearest listed wing can sit farther out than the nominal $w$ --- so
the return on capital at risk is bounded below by $-1$ exactly, and the
scoring asserts that bound. The hedge is asymmetric by design: on
long-volatility days the rule holds the plain package of
Equation~\eqref{eq:close_opt_settle}, whose loss is already capped at
the premium paid, and wings there would cap the tail the position
exists to own. Returns are quoted per dollar of body premium --- the
same denominator as the unhedged portfolio, so short and long days live in
one unit and the cost of the constraint reads directly against the
plain portfolio. The bounded capital-at-risk view is kept as the second
frame for the fraction-of-wealth exercise below, not as a headline: its
denominator is smallest exactly on rich-credit, high-volatility days.
Because the constraint is a venue fact rather than a modeling choice,
the questions asked of it are its price and the width and fraction to
choose under it, not whether the wings improve the portfolio.
\fi

\paragraph{Compounding at a fixed fraction of wealth.} The tables above
price one unit of premium per day and add P\&L up. A reinvesting
trader bets a fraction $f$ of current wealth instead, so wealth
compounds as $W_T=\prod_t(1+f\,q_t R_t)$ and the natural objective
is the annualized mean of $\log(1+f R'_t)$. Ruin is one bad day:
any $f \ge 1/|\min_t R'_t|$ puts wealth at or below zero on the worst
day, so the unhedged short portfolio, with a worst day near $-10$ premium
units, can only ever bet a small fraction of wealth no matter how
good its mean. The fraction is fixed at $f = 0.03$ --- three percent of
wealth per day, a chosen round number, the same on every day and for both
rules; nothing is estimated --- and every realized wealth factor is
asserted positive. Here $f$ is a share of wealth deployed as body premium,
the unit of the tables above. A defined-risk variant, which would bound the
worst day by construction, is parked and not reported here.

\paragraph{Scoring and control.} The unit of account is the expiration
day. For a portfolio of $n$ daily returns,
$t = \sqrt{n}\,\overline{R'}/\hat{s}(R')$, and the annualized Sharpe
ratio is $\overline{R'}/\hat{s}(R')\times\sqrt{252}$; the two carry the
same information at fixed $n$, and both are reported so the tables read
either way. The $\sqrt{252}$ is a per-\emph{trade}-day convention: one row
of the daily series is one trading period, and $252$ of them make a year.
This frame does not trade $252$ days a year. SPXW listed Monday, Wednesday
and Friday expirations before June~2022 and every session after, so over
2020--2024 it trades 200.3 days a year and reaches $252$ only from
mid-2022 on. A calendar-time annualization --- fill the non-expiration
sessions with a zero return and keep $\sqrt{252}$ --- multiplies every
annualized figure in this section by 0.892 and leaves every comparison
between rules and columns unchanged. The convention is stated rather than
switched, because the per-trade-day figure is the one a reader can
recompute from a table row's own mean and standard deviation. Fills are at
the 15:30 midpoint throughout, and no borrow or
market-impact model is applied. The one alternative fill that is reported
is the \emph{crossed spread}: both legs bought at the ask when the rule is
long and both sold at the bid when it is short, at entry only, since the
position is held into a cash settlement and pays no exit spread. It is
quoted beside the midpoint figure in
Section~\ref{sec:close_option_results} rather than replacing it, because
the midpoint figure is the one every other number in this section is
computed at.

The always-short control is the 15:30 end of the clock in
Section~\ref{sec:related}: \citet{HizmeriBevilacqua2025} locate the
predictive content of volatility uncertainty at 10:00~EST and gone by
midday; by the last bar the quoted implied already embeds a premium
over remaining realized variance. On the scaled comparison of
Equation~\eqref{eq:close_opt_ivscale}, implied variance exceeds the
forecast on 59 to 68\% of days depending on the column
(60\% for the headline), so the portfolio leans short more often
than long. What the forecast contributes at this hour is the
\emph{sign} of the day --- which side of the quoted premium to stand
on. Sizing beyond the sign was tested rather than assumed, and it is not
retained: scaling the position by the rank of $|s_t|$ within its own
trailing distribution raised no portfolio's Sharpe ratio, and that overlay
is parked rather than reported. The compounding exercise of
Section~\ref{sec:close_option_results} bets a fixed fraction of wealth and
reads nothing from the signal.
